\introduction
Dans un monde où les technologies de l’information évoluent rapidement, les entreprises de télécommunications doivent sans cesse moderniser leurs infrastructures pour répondre aux besoins croissants des clients. Pour cela, il est devenu indispensable de s’affranchir des systèmes anciens et fragmentés au profit d'un \textbf{inventaire unifié} capable de garantir la fiabilité des données et l'automatisation des services.\\
Ce rapport présente le travail que j’ai réalisé dans le cadre de mon stage de fin d’études au sein de l’entreprise \textbf{Inetum Offshore Maroc} pour le compte de son client \textbf{Operateur Telecom}. J’ai eu l’opportunité de participer à un projet technique d'envergure concernant la refonte du Système d'Information commercial, un programme visant à centraliser l'intégralité du patrimoine réseau et client dans la solution \textbf{Ni2 Unified Inventory}.\\
Afin de structurer ce rapport de manière claire et logique, j’ai choisi de regrouper mon travail en deux grandes parties. Cette organisation permet de mieux expliquer les missions que j’ai réalisées, en distinguant d'une part la \textbf{migration} des infrastructures techniques et d'autre part le défi de la \textbf{réconciliation} des données commerciales.\\
La première partie du projet concerne la migration du référentiel technique \textbf{GAIA}. Ce système gère l'ensemble du réseau cuivre et fibre de l'opérateur. Mon rôle a été d'analyser les structures de données existantes et de concevoir des scripts d'extraction \textbf{SQL} complexes. L'objectif était de reconstruire fidèlement la topologie du réseau dans le nouvel outil \textbf{Ni2}, tout en assurant une précision totale sur les ports, les équipements et les adresses de raccordement.\\
La deuxième partie du projet porte sur le volet \textbf{CRM}, qui représentait le défi majeur de ce stage. Il s'agissait d'une base de données massive sans documentation. J'ai dû mener un travail de \textbf{rétro-ingénierie} pour identifier les données clients critiques et les réconcilier avec l'inventaire technique. J'ai également mis en place un \textbf{pipeline de déploiement} industriel, incluant le build de l'application et l'automatisation de l'importation via des scripts Shell sur des serveurs distants, afin de garantir une intégration sécurisée et performante.